\documentclass[11pt,a4paper]{report}

% =============================================================================
% PACKAGES & IMPORTS
% =============================================================================
\usepackage[utf8]{inputenc}
\usepackage[french]{babel}
\usepackage[T1]{fontenc}
\usepackage{lmodern}
\usepackage[top=2.5cm,bottom=2.5cm,left=2.5cm,right=2.5cm]{geometry}
\usepackage{graphicx}
\usepackage{xcolor}
\usepackage{tikz}
\usetikzlibrary{shapes.geometric, arrows.meta, positioning, shadows, calc, fit}
\usepackage{amsmath,amssymb}
\usepackage{booktabs}
\usepackage{longtable}
\usepackage{array}
\usepackage{pifont}
\usepackage{fancyhdr}
\usepackage{titlesec}
\usepackage{enumitem}
\usepackage[most]{tcolorbox}
\usepackage{listings}
\usepackage{microtype}
\usepackage[pdftex,
    colorlinks=true,
    linkcolor=blue!70!black,
    citecolor=blue!70!black,
    urlcolor=blue!70!black,
    pdftitle={AutoCare AI - Cahier des Charges Version 3.0},
    pdfauthor={Equipe AutoCare AI},
    pdfsubject={Projet de Fin d'Etudes - Specifications Fonctionnelles et Techniques}]{hyperref}

% =============================================================================
% PALETTE DE COULEURS PROFESSIONNELLE
% =============================================================================
\definecolor{primaryBlue}{RGB}{20, 50, 110}
\definecolor{secondaryBlue}{RGB}{41, 128, 185}
\definecolor{accentBlue}{RGB}{171, 199, 255}
\definecolor{darkGray}{RGB}{40, 40, 40}
\definecolor{lightGray}{RGB}{245, 247, 250}
\definecolor{borderGray}{RGB}{210, 215, 225}
\definecolor{codeGreen}{RGB}{46, 139, 87}
\definecolor{codePurple}{RGB}{128, 0, 128}
\definecolor{successGreen}{RGB}{39, 174, 96}
\definecolor{warningOrange}{RGB}{230, 126, 34}

% =============================================================================
% CONFIGURATION DES STYLES DE TITRES ET EN-TÊTES
% =============================================================================
\titleformat{\chapter}[display]
  {\normalfont\bfseries\color{primaryBlue}}
  {\Large\color{secondaryBlue}Chapitre \thechapter}
  {1ex}
  {\Huge\vspace{-0.5ex}}
[\vspace{1ex}\titlerule]

\titleformat{\section}
  {\normalfont\Large\bfseries\color{primaryBlue}}
  {\thesection}{1em}{}

\titleformat{\subsection}
  {\normalfont\large\bfseries\color{secondaryBlue}}
  {\thesubsection}{1em}{}

\pagestyle{fancy}
\fancyhf{}
\fancyhead[L]{\small\color{gray}\textit{AutoCare AI — Cahier des Charges v3.0}}
\fancyhead[R]{\small\color{gray}\textit{Specifications Fonctionnelles \& Techniques}}
\fancyfoot[C]{\thepage}
\renewcommand{\headrulewidth}{0.4pt}
\renewcommand{\footrulewidth}{0.4pt}

% =============================================================================
% TCOLORBOX CONFIGURATION
% =============================================================================
\newtcolorbox{infoBox}[1][]{
  colback=lightGray,
  colframe=primaryBlue,
  fonttitle=\bfseries\color{white},
  coltitle=white,
  title=#1,
  arc=3mm,
  boxrule=1pt,
  left=10pt,right=10pt,top=8pt,bottom=8pt
}

\newtcolorbox{techBox}[1][]{
  colback=white,
  colframe=secondaryBlue,
  fonttitle=\bfseries\color{white},
  title=#1,
  arc=2mm,
  boxrule=0.8pt
}

% =============================================================================
% DEBUT DU DOCUMENT
% =============================================================================
\begin{document}

% -----------------------------------------------------------------------------
% PAGE DE GARDE
% -----------------------------------------------------------------------------
\begin{titlepage}
    \centering
    \begin{tikzpicture}[remember picture,overlay]
        \fill[primaryBlue] (current page.north west) rectangle ([yshift=-4cm]current page.north east);
        \fill[secondaryBlue] ([yshift=-4cm]current page.north west) rectangle ([yshift=-4.3cm]current page.north east);
    \end{tikzpicture}
    
    \vspace{0.5cm}
    {\Huge \bfseries \color{white} AutoCare AI} \\[0.4cm]
    {\Large \color{accentBlue} Cahier des Charges Fonctionnel et Technique} \\[2.5cm]
    
    \begin{tcolorbox}[colback=white,colframe=primaryBlue,arc=4mm,boxrule=1.5pt,halign=center]
        \vspace{0.3cm}
        {\Huge \bfseries \color{primaryBlue} CAHIER DES CHARGES} \\[0.3cm]
        {\Large \bfseries \color{secondaryBlue} Version 3.0} \\[0.2cm]
        {\small \color{darkGray} Projet de Fin d'Etudes (PFE)}
        \vspace{0.3cm}
    \end{tcolorbox}
    
    \vfill
    
    \begin{minipage}{0.48\textwidth}
        \begin{techBox}[Informations Generales]
            \textbf{Auteur :} Etudiant PFE \\
            \textbf{Projet :} AutoCare AI \\
            \textbf{Date :} Juillet 2026 \\
            \textbf{Statut :} Document de Ref\'erence v3.0
        \end{techBox}
    \end{minipage}
    \hfill
    \begin{minipage}{0.48\textwidth}
        \begin{techBox}[Technologies Cibles]
            \textbf{Backend :} Spring Boot 3 (Java 21) \\
            \textbf{Mobile :} React Native / Expo (TS) \\
            \textbf{Admin Web :} React / Tailwind CSS \\
            \textbf{Base de Donnees :} PostgreSQL \\
            \textbf{Intelligence Artificielle :} Google Gemini
        \end{techBox}
    \end{minipage}
    
    \vspace{1.5cm}
    
    \begin{center}
        \small \color{gray}
        Ce document definit les besoins, les fonctionnalites, les contraintes techniques et le perimetre du projet AutoCare AI dans le cadre d'un Projet de Fin d'Etudes.
    \end{center}
\end{titlepage}

% -----------------------------------------------------------------------------
% SOMMAIRE / TABLE DES MATIERES
% -----------------------------------------------------------------------------
\newpage
\tableofcontents
\newpage

% =============================================================================
% CHAPITRE 1 : INTRODUCTION ET GLOSSAIRE
% =============================================================================
\chapter{Introduction et Glossaire}

\section{Introduction}
Le present document constitue le cahier des charges fonctionnel et technique du projet \textbf{AutoCare AI}. Ce projet s'inscrit dans le cadre d'un Projet de Fin d'Etudes (PFE) et vise a concevoir et developper une solution logicielle pour optimiser et automatiser le suivi de maintenance des vehicules. A travers une application mobile destinee aux conducteurs et une console web dediee aux administrateurs, le systeme centralise les donnees techniques et modelise l'etat d'usure des pieces sous la forme d'un profil numerique de suivi, avec l'appui de l'intelligence artificielle pour la structuration initiale des donnees constructeurs.

\section{Glossaire}
\begin{description}
    \item[Jumeau Numerique (Digital Twin)] : Representation virtuelle et structuree d'un vehicule physique au sein de l'application, modelisant son etat actuel, ses specifications et son calendrier de maintenance.
    \item[Composant (Piece)] : Element mecanique ou consommable du vehicule (ex. huile moteur, plaquettes de frein, batterie) suivi individuellement en termes d'usure ou d'echeance.
    \item[Onboarding] : Processus d'inscription et de configuration initiale du profil d'un utilisateur ou de son vehicule dans l'application.
    \item[Filtre d'usure] : Algorithme ou regle permettant d'estimer la degradation ou le besoin de maintenance d'une piece en fonction du temps ecoule ou des kilometres parcourus.
\end{description}

\section{Acronymes}
\begin{itemize}[label=\textbullet]
    \item \textbf{API} : Application Programming Interface
    \item \textbf{JWT} : JSON Web Token
    \item \textbf{REST} : Representational State Transfer
    \item \textbf{UI / UX} : User Interface / User Experience
    \item \textbf{UML} : Unified Modeling Language
    \item \textbf{SaaS} : Software as a Service
    \item \textbf{PFE} : Projet de Fin d'Etudes
\end{itemize}

% =============================================================================
% CHAPITRE 2 : ANALYSE DU BESOIN ET PROBLEMATIQUE
% =============================================================================
\chapter{Analyse du Besoin et Problematique}

\section{Pourquoi ? (Le besoin)}
L'entretien regulier d'un vehicule est indispensable pour assurer la securite des passagers, reduire la consommation de carburant et eviter les pannes majeures. Pourtant, de nombreux conducteurs ne disposent pas d'un suivi clair des interventions necessaires ou effectuees. Les informations sont eparpillees dans des factures physiques, des carnets d'entretien papier souvent egarés, ou reposent sur des souvenirs approximatifs.

\section{Pour qui ? (Les cibles)}
La solution s'adresse a deux categories distinctes d'utilisateurs :
\begin{enumerate}
    \item \textbf{Les conducteurs particuliers (utilisateurs mobiles)} qui souhaitent simplifier et fiabiliser la gestion de leur(s) vehicule(s), conserver un historique d'entretien coherent et recevoir des alertes personnalisees.
    \item \textbf{Les gestionnaires de la plateforme (administrateurs web)} chargeux de superviser les comptes utilisateurs, de valider et d'enrichir le catalogue des marques et modeles de vehicules, et de suivre les files d'attente d'analyse de donnees.
\end{enumerate}

\section{Quel probleme ? (La problematique)}
La gestion classique de la maintenance automobile pose plusieurs problemes majeurs :
\begin{itemize}[label=\textbullet]
    \item \textbf{L'absence de proactivite} : Les conducteurs attendent souvent la panne ou l'allumage d'un voyant avant d'agir.
    \item \textbf{La perte de valeur a la revente} : L'incapacite a prouver le suivi regulier du vehicule deprecie sa valeur sur le marche de l'occasion.
    \item \textbf{La difficulte d'identification des pieces} : Les references constructeurs (OEM) et les specifications des fluides sont complexes pour les non-inities.
\end{itemize}

\section{Quelle solution ?}
\textbf{AutoCare AI} repond a ces defis en proposant une plateforme integree :
\begin{itemize}[label=\textbullet]
    \item \textbf{L'automatisation de la creation du profil vehicule} : A partir d'un assistant en ligne sur mobile, le systeme s'appuie sur un catalogue de reference (166+ marques, 1270+ modeles) et interroge un modele d'intelligence artificielle pour generer immediatement une cartographie de maintenance coherente (15+ composants standardises, intervalles et guides).
    \item \textbf{La centralisation documentaire} : Un espace dedie pour stocker et associer les guides constructeurs et factures d'entretien.
    \item \textbf{Un tableau de bord intuitif} : Permettant de visualiser en un coup d'oeil la sante mecanique du vehicule.
\end{itemize}

% =============================================================================
% CHAPITRE 3 : UTILISATEURS, ROLES ET CAS D'UTILISATION
% =============================================================================
\chapter{Utilisateurs, Roles et Cas d'Utilisation}

\section{Acteurs et Roles}
Le systeme distingue deux roles principaux disposant de droits d'acces specifiques :
\begin{itemize}[label=\textbullet]
    \item \textbf{Conducteur / Utilisateur Mobile} : Il gere son propre profil, ajoute ses vehicules, suit l'usure de ses pieces de rechange, declare ses interventions et telecharge ses documents.
    \item \textbf{Administrateur Web} : Il gere les utilisateurs et les vehicules, administre la base de reference des marques/modeles et suit l'execution des traitements de fonds du systeme.
\end{itemize}

\section{Diagramme des Cas d'Utilisation}
Le diagramme UML ci-dessous decrit les interactions des acteurs avec le systeme :

\begin{figure}[h!]
\centering
\begin{tikzpicture}[scale=0.8, every node/.style={transform shape}, >=Stealth]
    % Acteurs
    \node [circle, draw, minimum size=1cm, fill=lightGray] (user) at (-4, 2) {};
    \draw (user.south) -- ++(0,-1.2) coordinate (body) -- ++(-0.6,-1) -- ++(0.6,1) -- ++(0.6,-1);
    \draw (body) -- ++(0,0.8) coordinate (arms) -- ++(-0.8,0.2) -- ++(0.8,-0.2) -- ++(0.8,0.2);
    \node [below=1.5cm of user, align=center] {Utilisateur Mobile\\(Conducteur)};

    \node [circle, draw, minimum size=1cm, fill=lightGray] (admin) at (8, 2) {};
    \draw (admin.south) -- ++(0,-1.2) coordinate (body2) -- ++(-0.6,-1) -- ++(0.6,1) -- ++(0.6,-1);
    \draw (body2) -- ++(0,0.8) coordinate (arms2) -- ++(-0.8,0.2) -- ++(0.8,-0.2) -- ++(0.8,0.2);
    \node [below=1.5cm of admin, align=center] {Administrateur Web};

    % System boundary
    \draw[borderGray, thick] (-1.5, -4) rectangle (5.5, 7.5);
    \node[anchor=north west] at (-1.3, 7.3) {\textbf{Systeme AutoCare AI}};

    % Cas d'utilisation
    \node[draw, ellipse, fill=white, minimum width=2.5cm, align=center] (uc1) at (2, 6) {S'authentifier};
    \node[draw, ellipse, fill=white, minimum width=2.5cm, align=center] (uc2) at (2, 4) {Ajouter / Gerer\\un vehicule};
    \node[draw, ellipse, fill=white, minimum width=2.5cm, align=center] (uc3) at (2, 1.5) {Consulter le\\Jumeau Numerique};
    \node[draw, ellipse, fill=white, minimum width=2.5cm, align=center] (uc4) at (2, -1) {Declarer une\\intervention};
    \node[draw, ellipse, fill=white, minimum width=2.5cm, align=center] (uc5) at (2, -3) {Gerer le catalogue\\constructeur};

    % Relations User
    \draw[->] (-2.2, 1.5) -- (uc1);
    \draw[->] (-2.2, 1.3) -- (uc2);
    \draw[->] (-2.2, 1.1) -- (uc3);
    \draw[->] (-2.2, 0.9) -- (uc4);

    % Relations Admin
    \draw[->] (6.2, 1.5) -- (uc1);
    \draw[->] (6.2, 1.3) -- (uc2);
    \draw[->] (6.2, 0.9) -- (uc5);
\end{tikzpicture}
\caption{Diagramme UML de Cas d'Utilisation Global}
\label{fig:usecase}
\end{figure}

\section{User Stories Principales}
\begin{infoBox}[US-01 : Ajout rapide de vehicule]
    \textbf{En tant que} conducteur, \\
    \textbf{Je veux} pouvoir selectionner la marque, le modele et l'annee de mon vehicule via un formulaire guide, \\
    \textbf{Afin de} laisser le systeme generer automatiquement mon profil d'entretien numerique sans avoir a saisir manuellement toutes les pieces et preconisations.
\end{infoBox}

\begin{infoBox}[US-02 : Consultation de l'etat des pieces]
    \textbf{En tant que} conducteur, \\
    \textbf{Je veux} voir le pourcentage d'etat de mes bougies, huiles et plaquettes de frein, \\
    \textbf{Afin de} planifier mon prochain passage au garage avant que les pieces ne s'abiment severement.
\end{infoBox}

\begin{infoBox}[US-03 : Gestion administrative]
    \textbf{En tant qu'}administrateur, \\
    \textbf{Je veux} pouvoir editer ou ajouter des marques et modeles de vehicules dans la base de donnees globale, \\
    \textbf{Afin de} maintenir la base de donnees a jour avec les nouveaux modeles sortant sur le marche.
\end{infoBox}

% =============================================================================
% CHAPITRE 4 : SPECIFICATIONS FONCTIONNELLES
% =============================================================================
\chapter{Specifications Fonctionnelles}

\section{Matrice des Exigences Fonctionnelles (EF)}
Les exigences fonctionnelles majeures de la plateforme sont synthetisees dans le Tableau~\ref{tab:exigences_fonc}.

\begin{longtable}{@{}l p{6cm} c p{6cm}@{}}
\caption{Matrice des Exigences Fonctionnelles} \label{tab:exigences_fonc} \\
\toprule
\textbf{Code} & \textbf{Description de l'Exigence} & \textbf{Priorite} & \textbf{Module / Cible} \\ \midrule
\endhead
\toprule
\endfoot
\bottomrule
\endlastfoot
EF-AUTH-01 & Inscription et connexion de l'utilisateur avec validation d'adresse email. & Elevee & Securite (Mobile \& Web) \\
EF-AUTH-02 & Gestion de session sécurisée (Authentification double-jeton : Acces et Refresh). & Elevee & Securite (Mobile \& Web) \\
EF-VEH-01 & Assistant interactif d'onboarding de vehicules en 5 etapes clés. & Elevee & Gestion Vehicule (Mobile) \\
EF-VEH-02 & Recherche dynamique par marque, modele, annee et finitions constructeur. & Elevee & Catalogue (Mobile \& Web) \\
EF-TWIN-01 & Generation asynchrone du Jumeau Numerique incluant au moins 15 composants cles. & Elevee & Intelligence Artificielle \\
EF-TWIN-02 & Affichage de l'etat d'usure des pieces avec bareme de couleurs et barres de progression. & Elevee & Jumeau Numerique (Mobile) \\
EF-DOC-01 & Generation automatique de 5 documents PDF de reference constructeurs. & Moyenne & Gestion Documentaire (Mobile) \\
EF-INT-01 & Calendrier et fiches de periodicite pour 8 intervalles de maintenance standards. & Elevee & Intervalles (Mobile) \\
EF-ADM-01 & Consultation de l'activite de la plateforme (comptes, vehicules, etats des tâches IA). & Elevee & Backoffice (Web Admin) \\
EF-ADM-02 & Gestion des sessions actives et invalidation a distance par l'administrateur. & Moyenne & Backoffice (Web Admin) \\
\end{longtable}

\section{Regles de Gestion (RG)}
\begin{enumerate}[label=\textbf{RG-\arabic*}]
    \item \textbf{RG-Email-Unique} : Deux comptes utilisateurs ne peuvent pas posseder la meme adresse email.
    \item \textbf{RG-Status-Onboarding} : Un vehicule en cours de generation par l'IA doit posseder le statut transitoire \textit{PENDING} et ne peut pas etre consulte tant que sa cartographie technique n'est pas complete.
    \item \textbf{RG-Age-Vehicule} : L'annee de fabrication saisie pour un vehicule doit etre comprise entre 1900 et l'annee en cours majorée d'un an.
    \item \textbf{RG-Validite-Jeton} : Un jeton d'acces expire doit etre automatiquement renouvele a l'aide du jeton de rafraichissement si ce dernier est toujours valide en base de donnees.
\end{enumerate}

\section{Diagramme de Sequence : Generation du Jumeau Numerique}
Le diagramme ci-dessous illustre l'echange lors de la creation d'un vehicule :

\begin{figure}[h!]
\centering
\begin{tikzpicture}[scale=0.8, every node/.style={transform shape}, >=Stealth]
    % Lifelines
    \draw[thick] (0, 0) -- (0, -7) node[below] {Utilisateur};
    \draw[thick] (3, 0) -- (3, -7) node[below] {App Mobile};
    \draw[thick] (6, 0) -- (6, -7) node[below] {REST API};
    \draw[thick] (9, 0) -- (9, -7) node[below] {Background Job};
    \draw[thick] (12, 0) -- (12, -7) node[below] {Gemini API};

    % Steps
    \draw[->] (0, -1) -- (3, -1) node[midway, above] {1. Renseigne Vehicule};
    \draw[->] (3, -2) -- (6, -2) node[midway, above] {2. POST /vehicles};
    \draw[->] (6, -3) -- (9, -3) node[midway, above] {3. Soumet Job (Asynchrone)};
    \draw[->] (6, -3.5) -- (3, -3.5) node[midway, above] {4. Retourne Job ID};
    \draw[->] (9, -4.5) -- (12, -4.5) node[midway, above] {5. Analyse (Prompt structur\'e)};
    \draw[<-] (9, -5.5) -- (12, -5.5) node[midway, above] {6. Reponse JSON};
    \draw[->] (9, -6.5) -- (6, -6.5) node[midway, above] {7. Sauvegarde & Cloture Job};
\end{tikzpicture}
\caption{Diagramme de Sequence Simplifie de la Generation du Jumeau Numerique}
\label{fig:sequence}
\end{figure}

% =============================================================================
% CHAPITRE 5 : SPECIFICATIONS TECHNIQUES ET ARCHITECTURE
% =============================================================================
\chapter{Specifications Techniques et Architecture}

\section{Contraintes Techniques}
La plateforme doit respecter les architectures et technologies retenues pour le projet :
\begin{itemize}[label=\textbullet]
    \item \textbf{Backend} : Developpe en \textbf{Java 21} avec le framework \textbf{Spring Boot 3} et l'ecosysteme \textbf{Spring Security} pour la gestion de l'authentification et des roles.
    \item \textbf{Mobile} : Developpe en \textbf{React Native} et \textbf{TypeScript} sous l'environnement \textbf{Expo}.
    \item \textbf{Web Administrateur} : Construit en \textbf{React} avec le framework de style \textbf{Tailwind CSS}.
    \item \textbf{Base de Donnees} : \textbf{PostgreSQL} avec gestion de l'evolution du schema par des scripts de migration versionnes.
    \item \textbf{Intelligence Artificielle} : Int\'egration directe de l'API \textbf{Google Gemini} pour la generation de profils automobiles structures.
\end{itemize}

\section{Diagramme de Deploiement Physique}
Le diagramme de deploiement UML de la Figure~\ref{fig:deploy} illustre la distribution physique des differents composants du systeme :

\begin{figure}[h!]
\centering
\begin{tikzpicture}[scale=0.9, every node/.style={transform shape}, node distance=2.5cm, >=Stealth]
    % Node client Mobile
    \node[draw, rectangle, fill=lightGray!40, minimum width=3.5cm, minimum height=1.5cm, align=center] (mobileNode) {
        \textbf{Terminal Mobile} \\
        \small [Expo Client] \\
        React Native App
    };
    
    % Node client Web
    \node[draw, rectangle, fill=lightGray!40, minimum width=3.5cm, minimum height=1.5cm, align=center, below=of mobileNode] (webNode) {
        \textbf{Navigateur Client} \\
        \small [Web Admin Console] \\
        React / Tailwind
    };

    % Node serveur backend
    \node[draw, rectangle, fill=accentBlue!30, minimum width=4cm, minimum height=2.2cm, align=center, right=3cm of mobileNode, yshift=-1cm] (backendNode) {
        \textbf{Serveur d'Application} \\
        \small [API Backend Node] \\
        Spring Boot 3 (Java 21) \\
        Spring Security / Job Engine
    };

    % Node Base de donnees
    \node[draw, cylinder, fill=secondaryBlue!30, shape border rotate=90, minimum width=2.5cm, minimum height=2cm, align=center, right=3.5cm of backendNode] (dbNode) {
        \textbf{Serveur BD} \\
        PostgreSQL
    };

    % Cloud externalisé API Gemini
    \node[draw, cloud, cloud puffs=11, fill=primaryBlue!20, minimum width=3cm, minimum height=1.8cm, align=center, below=2cm of dbNode] (geminiNode) {
        \textbf{Service IA Externalis\'e} \\
        Google Gemini API
    };

    % Liaisons
    \draw[thick] (mobileNode) -- node[above, align=center] {\small REST/HTTPS \\ \small Port 8080} (backendNode);
    \draw[thick] (webNode) -- node[below, align=center] {\small REST/HTTPS \\ \small Port 8080} (backendNode);
    \draw[thick] (backendNode) -- node[above] {\small JDBC / SQL} (dbNode);
    \draw[thick] (backendNode) -- node[left, align=center] {\small HTTP REST \\ \small JSON} (geminiNode);

\end{tikzpicture}
\caption{Diagramme de Deploiement du Systeme AutoCare AI}
\label{fig:deploy}
\end{figure}

\section{Data Flow Diagram (DFD)}
Le flux des donnees au sein du systeme s'articule comme suit :
\begin{itemize}[label=\textbullet]
    \item \textbf{Flux d'inscription} : Les informations d'authentification transitent depuis le client mobile ou web vers le controleur de securite, qui chiffre le mot de passe et l'enregistre en base de donnees. Un jeton JWT contenant les roles est retourne.
    \item \textbf{Flux de creation de vehicule} : Les criteres saisis par l'utilisateur sont transmis au service backend, qui verifie si les references de marques et modeles existent dans la base locale. Le service envoie ensuite une requete formateee au modele d'IA. La reponse de ce dernier est verifiee syntaxiquement avant d'etre parseee et stockee dans les tables correspondantes.
\end{itemize}

% =============================================================================
% CHAPITRE 6 : MODELE DE DONNEES ET BASE DE DONNEES
% =============================================================================
\chapter{Modele de Donnees et Base de Donnees}

\section{Diagramme de Classes Simplifie}
Le modele conceptuel de donnees est illustre par le diagramme UML de la Figure~\ref{fig:class}.

\begin{figure}[h!]
\centering
\begin{tikzpicture}[scale=0.9, every node/.style={transform shape}, node distance=2cm, >=Stealth]
    % Classe User
    \node[draw, rectangle, fill=lightGray!40, minimum width=4.5cm, minimum height=2.2cm, align=left] (userClass) {
        \textbf{User} \\
        \hline
        - id : Long \\
        - email : String \\
        - passwordHash : String \\
        - status : String
    };

    % Classe Vehicle
    \node[draw, rectangle, fill=accentBlue!20, minimum width=4.5cm, minimum height=2.2cm, align=left, right=3.5cm of userClass] (vehicleClass) {
        \textbf{Vehicle} \\
        \hline
        - id : Long \\
        - licensePlate : String \\
        - purchaseCondition : String \\
        - purchaseMileage : Integer \\
        - currentMileage : Integer
    };

    % Classe Component
    \node[draw, rectangle, fill=secondaryBlue!20, minimum width=4.5cm, minimum height=2.2cm, align=left, below=3cm of vehicleClass] (compClass) {
        \textbf{Component} \\
        \hline
        - id : Long \\
        - name : String \\
        - category : String \\
        - healthScore : Integer \\
        - partNumber : String
    };

    % Relations
    \draw[thick] (userClass) -- node[above, near start] {1} node[above, near end] {0..*} (vehicleClass);
    \draw[thick] (vehicleClass) -- node[left, near start] {1} node[left, near end] {0..*} (compClass);

\end{tikzpicture}
\caption{Diagramme de Classes Conceptuel (Entites Principales)}
\label{fig:class}
\end{figure}

\section{Schema Relationnel et Migrations}
Le schema de base de donnees est gere de maniere incrementale a l'aide de scripts Flyway. Les grandes etapes de l'evolution du schema sont :
\begin{itemize}[label=\textbullet]
    \item \textbf{Gestion des comptes} : Creation des tables des utilisateurs, rôles, sessions actives et historique de securite.
    \item \textbf{Gestion des vehicules} : Creation de la table du catalogue des marques et modeles, de la table des vehicules lies aux conducteurs et de la table des jobs d'arriere-plan.
    \item \textbf{Modellisation du Jumeau Numerique} : Table des composants (nom, categorie, numero de reference standard constructeur) et table des intervalles preconises.
\end{itemize}

% =============================================================================
% CHAPITRE 7 : EXIGENCES NON-FONCTIONNELLES ET SECURITE
% =============================================================================
\chapter{Exigences Non-Fonctionnelles et Securite}

\section{Exigences Non-Fonctionnelles (ENF)}
Les exigences non-fonctionnelles definissent les standards de qualite du systeme. Elles sont detaillees dans le Tableau~\ref{tab:exigences_non_fonc}.

\begin{table}[h!]
\centering
\caption{Matrice des Exigences Non-Fonctionnelles}
\label{tab:exigences_non_fonc}
\vspace{0.2cm}
\begin{tabular}{@{}l p{6cm} c p{6cm}@{}}
\toprule
\textbf{Code} & \textbf{Description de l'Exigence} & \textbf{Priorite} & \textbf{Cible} \\ \midrule
ENF-PERF-01 & Temps de reponse moyen des APIs REST inferieur a 200 ms hors requete vers l'IA. & Elevee & Performance du Backend \\
ENF-DISP-01 & Taux de disponibilite de la plateforme superieur a 99\% en conditions normales. & Elevee & Infrastructure \\
ENF-SEC-01 & Chiffrement systematique des mots de passe des utilisateurs en base de donnees. & Elevee & Securite des Donnees \\
ENF-SEC-02 & Validation de la validite des jetons JWT a chaque requete privee. & Elevee & Securite des APIs \\
ENF-UTIL-01 & Interface mobile ergonomique et responsive offrant un guidage clair de l'utilisateur. & Elevee & Ergonomie (UI/UX) \\
ENF-CONF-01 & Isolation complete des donnees : un conducteur ne peut voir que ses vehicules. & Elevee & Securite des Donnees \\
\bottomrule
\end{tabular}
\end{table}

\section{Securite et Confidentialite}
\begin{itemize}[label=\textbullet]
    \item \textbf{Chiffrement au repos} : Toutes les informations sensibles (telles que les mots de passe) sont cryptees en utilisant l'algorithme standard BCrypt avec un facteur de securite adapte.
    \item \textbf{Securisation des echanges} : Les flux d'API utilisent exclusivement le protocole securise HTTPS. Les en-tetes CORS (Cross-Origin Resource Sharing) sont configures pour restreindre les acces non autorises.
\end{itemize}

% =============================================================================
% CHAPITRE 8 : PLANIFICATION ET CRITERES D'ACCEPTATION
% =============================================================================
\chapter{Planification et Criteres d'Acceptation}

\section{Livrables Attendus}
Le projet doit fournir les livrables suivants a l'issue du stage :
\begin{enumerate}
    \item \textbf{Le code source du backend} Spring Boot 3 documente par une specification OpenAPI/Swagger.
    \item \textbf{Le code source de l'application mobile} React Native (Expo) compilable et installable pour test.
    \item \textbf{Le code source de la console web} d'administration React.
    \item \textbf{La base de donnees PostgreSQL} initialisee avec son historique de scripts de migration Flyway.
    \item \textbf{Les documentations associees} : Rapport de PFE et manuel d'installation technique.
\end{enumerate}

\section{Criteres d'Acceptation pour Validation (Recette)}
La validation finale repose sur le succes des tests suivants :
\begin{itemize}[label=\textbullet]
    \item[\checkmark] L'utilisateur peut s'inscrire, verifier son compte par courriel et se connecter de maniere securisee.
    \item[\checkmark] L'assistant d'onboarding mobile complete les 5 etapes sans erreur de navigation et cree un vehicule.
    \item[\checkmark] La creation du vehicule declare un job d'arriere-plan, interroge l'IA et insere en base un jumeau numerique complet (composants cles avec specifications, intervals de vidange, documents associes).
    \item[\checkmark] L'administrateur de la console web peut voir les statistiques et invalider les sessions des utilisateurs connectes.
\end{itemize}

\section{Analyse des Risques}
\begin{table}[h!]
\centering
\caption{Matrice d'Analyse des Risques du Projet}
\label{tab:risques}
\vspace{0.2cm}
\begin{tabular}{@{}p{3.5cm} c p{4cm} p{6cm}@{}}
\toprule
\textbf{Risque} & \textbf{Severite} & \textbf{Impact} & \textbf{Plan d'Attenuation} \\ \midrule
Temps de latence de l'API IA & Moyenne & Blocage de l'interface client & Dechargement de l'appel dans une file de jobs asynchrones avec indicateur de chargement. \\
Indisponibilite du service IA & Elevee & Impossibilite de creer un vehicule & Mise en place d'un profil de donnees d'entretien par defaut et pre-enregistre. \\
Vol de donnees de session & Elevee & Acces non autorise aux profils & Utilisation de jetons a duree de vie courte et stockage crypte. \\
\bottomrule
\end{tabular}
\end{table}

\end{document}
